% Vorlage für WABs der Provadis Hochschule
% basiert auf "Universität Ulm Praktikumsbericht Vorlage"
% von Max Sch.
% CC BY 4.0
%
% Titelseite.tex
%
% Hier wird die Titelseite erzeugt.
% Pflichtangaben laut Anforderungen (Formatierung 1a-f): Art der Arbeit, Titel,
% Fach (nur Praxisberichte), Name, Matrikelnummer, Studiengang mit Studiengruppe,
% E-Mail, Telefon (optional), Abgabedatum, Referent:in (+ Korreferent:in bei BT/MT).
% Alle Werte werden in Berichtinfos.tex gesetzt.
\begin{titlepage}

    % Logo Provadis Hochschule und ggf. Logo Arbeitgeber
    \includegraphics[height=2.06cm]{Bilder/Deckblatt/provadis-hochschule.pdf}
    \hfill
    % optionales Unternehmenslogo
    %\includegraphics[height=2.06cm]{Bilder/Deckblatt/Infraserv_Höchst_logo.pdf}

    \vspace*{1cm}

    \begin{singlespace}
        \begin{center}
            \large
            \artderarbeit

            \vspace*{2cm}

            \Large

            \textbf{\titel}

            \vspace*{2.5cm}

            \normalsize
            Zur Veranstaltung / zur Erlangung des akademischen Grades\\
            ’Bachelor of Science’ B.Sc. / ’Master of Science’ M.Sc.\\
            im Studiengang ’XXX’
            \ifdefempty{\fach}{}{%
                im Fach \enquote{\fach}\\
            }%
            \ifdefempty{\abschluss}{}{%
                zur Erlangung des akademischen Grades\\
                \abschluss\\
            }

            \vspace*{1.5cm}

            vorgelegt dem Fachbereich \fakultaet\ der\\
            Provadis School of International Management and Technology\\
            von

            \vspace*{1cm}

            \fullname \\
            Matrikelnummer: \matnr \\
            Studiengang: \studiengang, Studiengruppe: \studiengruppe \\
            \email
            \ifdefempty{\telMobil}{}{\\ \telMobil}
        \end{center}
    \end{singlespace}

    \normalsize
    \vfill % variabler vertikaler Abstand
    \ifdefempty{\refSecond}%
        {\def\korefzeile{}}%
        {\def\korefzeile{Korreferent:in: & \refSecond \\[0.25cm]}}%
    \begin{tabular}{@{}ll}
        Referent:in: & \refFirst \\[0.25cm]
        \korefzeile
        Abgabedatum: & \abgabedatum \\
    \end{tabular}

\end{titlepage}

% Ende der Datei
